\documentclass[10pt,a4paper]{amsart}
\usepackage[margin=19mm]{geometry}
\usepackage{amsmath,amssymb,mathtools}
\usepackage[scr=rsfs,cal=euler]{mathalfa}
\usepackage{xfrac}
\usepackage[unicode,hidelinks,bookmarks=false]{hyperref}
\newcommand{\ie}{i.e.\ }
\newcommand{\cf}{cf.\ }
\newcommand{\resp}{respectively\ }
\newcommand{\Subsection}{\subsection}
\providecommand{\nobreakdash}{\nobreak\hbox{-}}
\setlength{\emergencystretch}{2em}
\multlinegap0pt
\usepackage{amssymb}
\usepackage{xcolor}
\newtheorem{theorem}{Theorem}
\newtheorem{lemma}{Lemma}
\newtheorem{proposition}{Proposition}
\usepackage{cleveref}
\crefname{theorem}{Theorem}{Theorems}
\crefname{lemma}{Lemma}{Lemmas}
\crefname{proposition}{Proposition}{Propositions}
\newcommand{\E}{\mathbb E}
\newcommand{\Ecbinf}{\mathcal E^{\mathrm{cb},\infty}}
\newcommand{\Hh}{\mathcal H}
\newcommand{\R}{\mathbb R}
\newcommand{\T}{\mathcal T}
\newcommand{\X}{\mathcal X}
\newcommand{\Z}{\mathcal Z}
\title{On the optimality of coin-betting for mean estimation}
\author{Eugenio Clerico}
\date{Source: arXiv:2412.02640v4 (CC BY 4.0)}
\begin{document}
\maketitle

\noindent\emph{Source and licence.} On the optimality of coin-betting for mean estimation. Version arXiv:2412.02640v4 (CC BY 4.0), distributed by arXiv under the Creative Commons Attribution 4.0 International licence, http://creativecommons.org/licenses/by/4.0/, as recorded in the arXiv licence metadata for this exact version and shown on the abstract page https://arxiv.org/abs/2412.02640v4. Full licence notice: \texttt{licenses/arxiv-2412.02640-CC-BY-4.0.txt}.\par\smallskip

\noindent\textit{Editorial scope.} Counted unit: the article's theorem thm:proc (source0.tex lines 321--323). The article's complete original proof of it is delivered with it (source0.tex lines 462--478, one \texttt{proof} environment of 17 lines, which the article places away from the statement). No mathematics is written, repaired or re-derived: the statement and the proof are the article's own lines, in the article's own notation, and no retained line is substituted or re-flowed.  Retained uncounted as the source-local context the proof needs: lines 430--437; lines 440--444; lines 521--521; lines 534--536; lines 550--552.  Nothing was omitted from any retained span, so the package carries no omission record.  No retained line contains a citation, so the wrapper carries no bibliography and no bibliography entry.  No retained line contains a figure environment, an image-inclusion command or drawing code, so no image is delivered and the package declares no asset dependency.  Editorial formatting only, in addition: an AMS \texttt{amsart} preamble that restates the article's own declarations and macros used by the retained lines, namely the environments \texttt{theorem}, \texttt{lemma}, \texttt{proposition} on separate counters, together with the standard packages those lines need.  The retained environments print in this document as Theorem 1, Lemma 1, Proposition 1 in order of appearance, whereas the article prints the same items with the numbering of its own full document.  The complete native source of the article as distributed in the arXiv e-print for this exact version is bundled unchanged as \texttt{sources/169551/source0.tex}.
\begin{lemma}\label{lemma:trees}
Let $f = (f_s)_{s\geq 1}$ denote a sequence of bounded Borel functions, with $f_s:\X^s\to\R$. Fix $t\geq 1$ and, for any $x^t\in\X^t$ and $s\geq0$, let $f_s^{x^t}:\X^s\to\R$ be defined via $f_s^{x^t}(y^s) = f_{t+s}(x^t,y^s)$. Define the functions $u$ and $l$, from $\X^t\to\R$, as
$$
    u(x^t) = \sup_{Q\in\bar\Hh^\infty_\mu}\sup_{\tau\in\T}\E_Q[f_\tau^{x^t}]\quad\text{ and }\quad
    l(x^t) = \inf_{Q\in\bar\Hh^\infty_\mu}\inf_{\tau\in\T}\E_Q[f_\tau^{x^t}]\,.
$$
Then, $u$ is upper semi-analytic and $l$ is lower semi-analytic.
\end{lemma}

\begin{lemma}\label{lemma:TQx}
    Let $\mu\in(0,1)$ and $E$ be an e-process for $\bar\Hh^\infty_\mu$. Then there is $\lambda_1\in I_\mu$ such that, for any $\tau\in\T$, any $Q\in\bar\Hh_\mu^\infty$, and any $x\in\X$, we have that
	$$\E_Q[E^x_\tau]\leq E_{\lambda_1}(x) = 1 + \lambda_1 (x-\mu)\,,$$
    where $E^x = (E^x_t)_{t\geq 0}$ is the sequence of Borel functions $E_t^x:\X^{t}\to I_\mu$, defined via $E_t^x(y^{t}) = E_{t+1}(x, y^t)$, for $t\geq 1$ and $y^t\in\X^t$, and $E_0^x = E_1(x)$.
\end{lemma}

\subsection{A functional variant of Lusin's separation theorem}\label{app:lusin}

\begin{lemma}\label{lemma:supinf}
	Let $\Z\subseteq\R^d$ and $\Z'\subseteq\R^{d'}$ be Borel sets. Let $f:\Z\times\Z'\to\overline\R$ be an upper semi-analytic function. Then $u:\Z\to\overline\R$ defined as $u(z) = \sup_{z'\in\Z'}f(z,z')$ is upper semi-analytic. Similarly, let $l:\Z\to\overline\R$ be given by $l(z) = \inf_{z'\in\Z'}f(z,z')$. Then, $l$ is lower semi-analytic.
\end{lemma}

\begin{proposition}\label{prop:ufl}
    Let $\Z\subseteq\R^d$ be a Borel set. Let $u:\Z\to\R$ be an upper semi-analytic function bounded from below and $l:\Z\to\R$ a lower semi-analytic function bounded from above. If $u\preceq l$, there exists a Borel function $b:\Z\to\R$ such that $u\preceq b\preceq l$. 
\end{proposition}

\begin{theorem}\label{thm:proc}
    For any $\mu\in(0,1)$ the coin-betting e-process class is the optimal e-process class for $\Hh_\mu^\infty$.
\end{theorem}

\begin{proof}[of \Cref{thm:proc}]
	First, let us show that all the e-processes in $\Ecbinf_\mu$ are maximal (for $\Hh_\mu^\infty$). Assume that this was not the case. Then there is a $\lambda^\infty\in\Lambda_\mu^\infty$ and an e-process (for $\Hh_\mu^\infty$) $E\succeq E^{\lambda^\infty}$ such that, for some $t\geq t$, $E_t\succ E_{\lambda^t}$ (clearly this cannot happen for $t=0$, as we must have $E_0\leq 1$ for every e-process). So, there is $\hat x^t\in\X^t$ such that $E_t(\hat x^t)>E_{\lambda^t}(\hat x^t)$. It is not hard to see that there is a $Q\in\Hh_\mu^\infty$ that puts non-zero mass on the set $\{x^\infty\in\X^\infty\,:\,x^t = \hat x^t\}$. We can consider the constant stopping time $\tau=t$. Then, we have that $0< E_t(\hat x^t) - E_{\lambda^t}(\hat x^t) \leq \E_Q[E_t - E_{\lambda^t}] = \E_Q[E_\tau] - 1 \leq  0$,  a contradiction.	

	Hence, we are left with showing that $\Ecbinf_\mu$ is a majorising e-process class for $\Hh^\infty_\mu$. Since every e-process in for $\Hh^\infty_\mu$ is also an e-process for $\bar\Hh^\infty_\mu$, it is sufficient to show that $\Ecbinf_\mu$ is a majorising e-process class for $\bar\Hh^\infty_\mu$. Fix an e-process $E$ (for $\bar\Hh^\infty_\mu$). We introduce the following notation. For $t\geq 1$ and $x^t\in\X^t$, we denote as $E^{x^t} = (E^{x^t}_s)_{s\geq 0}$ the sequence of non-negative functions $E^{x^t}_s:\X^{s-1}\to I_\mu$, defined via $E_s^{x^t}(y^s) = E_{t+s}(x^t, y^s)$, for any $y^s\in\X^s$ and $s\geq 1$, and $E_0^{x^t} = E_t(x^t)$. In what follows, we say that a $T$-tuple $\lambda^{T}$ of Borel functions $\lambda_i:\X^{i-1}\mapsto I_\mu$, \emph{dominates $E$ at level $T$} if, for all $Q\in\bar\Hh^\infty_\mu$, for any $t\in[1:T]$, for all $\tau\in \T$ and $x^t\in\X^t$, we have that $$\E_Q[E^{x^t}_\tau]\leq E_{\lambda^t}(x^t)\,.$$ We will show that there is a $\lambda^\infty\in\Lambda_\mu^\infty$ such that, for any $T\geq 1$, $\lambda^T$ dominates $E$.

    We construct this sequence progressively. By \Cref{lemma:TQx}, we know that there is $\lambda_1\in I_\mu$ that dominates $E$ at level $1$. Now, say that, for some $T\geq 2$, there is a $\lambda^{T-1}\in\Lambda_\mu^{T-1}$ that dominates $E$ at level $T-1$. 
Let us show that this imply the existence of a Borel $\lambda_T:\X^{T-1}\to I_\mu$, such that $\lambda^T=(\lambda^{T-1},\lambda_T)$ dominates $E$ at level $T$. Define $\mathcal S=\{x^{T-1}\in\X^{T-1}\,:\,E_{\lambda^{T-1}}(x^{T-1})\neq 0\}$. For any $x^{T-1}\in\mathcal S$, define the sequence of Borel functions $\tilde E^{x^{T-1}}  = (\tilde E^{x^{T-1}}_t)_{t\geq 0}$ as follows. $\tilde E_0^{x^{T-1}} = 1$, and $\tilde E_t^{x^{T-1}}(y^t) = E_{T-1+t}(x^{T-1}, y^t)/E_{\lambda^{T-1}}(x^{T-1})$ for $t\geq 1$ and $y^t\in\X^t$. By construction, $\tilde E^{x^{T-1}}$ is an e-process. In particular, by \Cref{lemma:TQx}, there must be $\tilde\lambda_1^{x^{T-1}}\in I_\mu$ (which of course might depend on $x^{T-1}$ in a non-Borel way) dominating $\tilde E^{x^{T-1}}$ at level $1$. So, for any $x^{T-1}\in\mathcal S$ and $y\in\X$,
$$\sup_{Q\in\bar\Hh^\infty_\mu}\sup_{\tau\in\T}\E_Q[  E_\tau^{(x^{T-1},y)}] \leq  E_{\lambda^{T-1}}(x^{T-1}) \big(1 +  \tilde\lambda_1^{x^{T-1}}(y-\mu)\big)\,.$$ 
We now define the mappings $u$ and $l$ on $\mathcal S$ as 
\begin{align*}
	u(x^{T-1}) &=\sup_{\tau\in\T}\sup_{Q\in\bar\Hh^\infty_\mu}\sup_{y\in\X\cap(\mu,1]}\frac{1}{y-\mu}\left(\frac{\E_Q[E^{(x^{T-1},y)}_\tau]}{E_{\lambda^{T-1}}(x^{T-1})}-1\right)\,;\\
   	l(x^{T-1}) &= \inf_{\tau\in\T}\inf_{Q\in\bar\Hh^\infty_\mu}\inf_{y\in\X\cap[0,\mu)}\frac{1}{\mu-y}\left(\frac{\E_Q[E^{(x^{T-1},y)}_\tau]}{E_{\lambda^{T-1}}(x^{T-1})}-1\right)\,.
\end{align*}
It follows from \Cref{lemma:trees} that $(x^{T-1},y)\mapsto \sup_{\tau\in\T}\sup_{Q\in\bar\Hh^\infty_\mu}\frac{1}{y-\mu}\big(\frac{\E_Q[E^{(x^{T-1},y)}_\tau]}{E_{\lambda^{T-1}}(x^{T-1})}-1\big)$ is upper semi-analytic (note that $\mathcal S$ is a Borel set, so the domain restriction does not cause any issue). In particular,  $u$ is also upper semi-analytic on $\mathcal S$ (see \Cref{lemma:supinf} in \ref{app:lusin}). By an analogous argument, $l$ is lower semi-analytic on $\mathcal S$. Moreover, $$(\mu-1)^{-1}\leq u(x^{T-1})\leq\tilde\lambda_1^{x^{T-1}}\leq l(x^{T-1})\leq\mu^{-1}$$ for all $x^{T-1}\in\mathcal S$, and so in particular by \Cref{prop:ufl} (see \ref{app:lusin}) there is a Borel function $\lambda_{T}:\mathcal S\to I_\mu$ such that $l(x^{T-1})\geq \lambda_{T}(x^{T-1}) \geq u(x^{T-1})$ for all $x^{T-1}\in\mathcal S$. We can extend $\lambda_{T}$ to the whole $\X^{T-1}$, by setting $\lambda_{T}(x^{T-1}) = 0$, if $x^{T-1}\in \X^{T-1}\setminus \mathcal S$. Noting that $x^{T-1}\in\X^{T-1}\setminus \mathcal S$ implies that $E_{T-1+t}(x^{T-1}, y^t)=0$ for any $t\geq 1$ and $y^t\in\X^{t}$,\footnote{Indeed, for every $t$ the non-negativity of $E_{T-1+t}$ implies that this must be true for $Q$-almost every $y^t$, for every $Q\in\bar\Hh^\infty_\mu$, and so for every $y^t\in\X^t$.} one can easily verify that $\lambda^{T}$ dominates $E$ at level $T$. 

So, we can construct iteratively a sequence $\lambda^\infty\in\Lambda_\mu^\infty$ such that, for each $T\geq 1$, $\lambda^T$ dominates $E$ at level $T$. It follows immediately that $E^{\lambda^\infty}\succeq E$, and so $\Ecbinf_\mu$ is a majorising e-process class for $\bar\Hh_\mu^\infty$. \qed
\end{proof}

\end{document}
